\section{First circuit: from source to result}

The shortest useful learning path is to prepare one qubit, transform it,
measure it, and declare it as an output:

\textbf{Complete runnable protocol}
\begin{lstlisting}
MAIN-PROCESS FirstCircuit
CREATETOKEN -I Q
SET Q 0p
H -I Q -O Q
MEASURE -I Q
RETURNVALS Q
\end{lstlisting}

\subsection{Read the program line by line}

\begin{enumerate}
  \item \code{MAIN-PROCESS FirstCircuit} gives the protocol a name.
  \item \code{CREATETOKEN -I Q} reserves a logical wire named Q.
  \item \code{SET Q 0p} prepares Q in the basis state $\ket0$.
  \item \code{H -I Q -O Q} applies Hadamard to that same wire, producing
  $\ket+=(\ket0+\ket1)/\sqrt2$.
  \item \code{MEASURE -I Q} samples 0 or 1 and collapses the state.
  \item \code{RETURNVALS Q} declares Q as the first public result.
\end{enumerate}

\begin{center}
\begin{tikzpicture}[scale=.95]
  \draw[thick] (0,0)--(7.7,0);
  \draw[double distance=1.4pt,thick] (0,-.85)--(7.7,-.85);
  \node[left] at (0,0) {$\ket0$ Q};
  \node[left] at (0,-.85) {c};
  \node[left,font=\scriptsize] at (0.35,-.55) {/1};
  \node[draw=QpuBlue,fill=QpuBlue!8,minimum width=.7cm,minimum height=.6cm,font=\bfseries] at (2,0) {H};
  \node[above] at (3.6,0) {$\ket+$};
  \QpuMeter{5.2}{0}
  \draw[double distance=1.4pt,thick,-{Latex[length=1.6mm]}] (5.2,0)--(5.2,-.85);
  \node[right,font=\scriptsize] at (5.2,-.62) {0};
  \node[above] at (6.6,-.85) {0 or 1};
\end{tikzpicture}
\end{center}

\subsection{What one run means}

One execution gives one random-looking sample. For this circuit, theory predicts
equal probability for 0 and 1. Getting 0 once does not mean H failed, and
getting 1 once does not prove a bias. Repeat the circuit to compare observed
frequencies with the predicted distribution.

\subsection{Adding user-controlled input}

Move a wire into PARAMS when callers or truth-table rows should choose its
start state:

\textbf{Complete runnable protocol}
\begin{lstlisting}
PARAMS: Input:state

MAIN-PROCESS Inverter
X -I Input -O Input
MEASURE -I Input
RETURNVALS Input
\end{lstlisting}
The matching deterministic truth table is:
\begin{center}
\begin{tabular}{c|c}\toprule Input & Output\\\midrule
\code{0p}&\code{1p}\\
\code{1p}&\code{0p}\\\bottomrule
\end{tabular}
\end{center}
The rest of this manual explains every symbol in these examples before moving
to multi-qubit gates, child processes, and full truth-table files.

